% Applied Software Development Lab -- Lab Handout 4: Fun and games with the serial interface
% Copyright (c) Michael Lehn. Licensed under CC BY-SA 4.0.
%
% Region coordinates (bp from the top of the page) follow the page layout of
% the original German handout "Laborblatt 4".
\documentclass{aslab}
\sethandout{4}{Fun and games with the serial interface}
\graphicspath{{../figures/handout04/}{../figures/shared/}}

\makeatletter
%% Lists used in this handout --------------------------------------------
%% learning objectives: blue bold dashes, text at +12bp
\newlist{bluedash}{itemize}{1}
\setlist[bluedash]{label={\color{aslabblue}\bfseries\aslab@dash},
  leftmargin=12bp,labelsep=*,labelindent=0pt,align=left}
%% "A)" lists with the text at +18bp
\newlist{alphatasks}{enumerate}{1}
\setlist[alphatasks]{label={\Alph*)},leftmargin=18bp,labelsep=*,align=left,
  labelindent=0pt}
%% plain "1." lists (regular weight)
\newlist{numlist}{enumerate}{1}
\setlist[numlist]{label={\arabic*.},leftmargin=18bp,labelsep=*,align=left,
  labelindent=0pt}

%% Program listings with syntax colors -----------------------------------
%% Inside asmlisting, every space is a fixed 0.6em and every line break
%% ends a line, so the columns of the original listings are preserved
%% (an empty listing line is written as \mbox{}). The optional argument is the font size in bp (default 11); the line
%% pitch is the one of the surrounding region.
\newcommand*\aN[1]{\textcolor{codeorange}{#1}}% line numbers, keywords
\newcommand*\aI[1]{\textcolor{codegreen}{#1}}% types
\newcommand*\aK[1]{\textcolor{codered}{#1}}% literals
\newcommand*\aR[1]{\textcolor{codecyan}{#1}}% functions, objects
\newcommand*\aC[1]{\textcolor{codepurple}{#1}}% comments
\newcommand*\aM[1]{\textcolor{codemagenta}{#1}}% Python import
\newcommand*\asm@space{\leavevmode\hbox to 0.6em{}}
{\catcode`\ =\active\catcode`\^^M=\active%
\gdef\asm@setup{\catcode`\ =\active\def {\asm@space}%
\catcode`\^^M=\active\def^^M{\par}}}%
\newenvironment{asmlisting}[1][11]{\par
  \edef\asm@bs{\the\baselineskip}\fontsize{#1bp}{\asm@bs}\selectfont
  \ttfamily\asm@setup}{\par}

%% Hyperlink in the heading blue, as in this handout of the original
\newcommand*\bluelink[2]{\href{#1}{\textcolor{aslabblue}{#2}}}

%% Keep the baseline pitch fixed even next to tall inline code boxes.
%% Regions marked `free' below follow a line at a distance of less than
%% 12.3bp, so only the gray inline-code backgrounds touch (as in the
%% original); their text does not overlap.
\lineskiplimit=-\maxdimen
\makeatother

\begin{document}
%% ===================================================================== 1 ==
\maketitleblock

\begin{region}{124.0}
\heading{Learning objectives}
\end{region}

\begin{region}[pitch=12.1]{138.0}
\begin{bluedash}
\item Understanding serial communication between the microcontroller and
  the computer.
\item Applying the basics of the Python programming language.
\item A little bit of \code{git}
\end{bluedash}
\end{region}

\begin{region}{188.8}
\heading{Step 1: Forking the repository}
\end{region}

\begin{region}[pitch=13.25]{217.4}
\textbf{Goal:} From now on, we will develop a software project step by step.
We will proceed one step at a time -- and things will surely go wrong now
and then. But that is exactly what we have git for -- a tool for version
control.
\end{region}

\begin{region}{270.2}
\begin{tasks}
\item Open the following URL in your web browser:
  \bluelink{http://gitlab.uni-ulm.de/mlehn/asp-ulm}{http://gitlab.uni-ulm.de/mlehn/asp-ulm}
\end{tasks}
\end{region}

\begin{region}{283.4}
\begin{tasks}[start=2]
\item Log in with your kiz account.
\end{tasks}
\end{region}

\begin{region}{297.5}
\begin{tasks}[start=3]
\item Click \textbf{Fork} at the top right.
\end{tasks}
\end{region}

\begin{region}[pitch=14.3]{311.7}
\begin{tasks}[start=4]
\item Under \textbf{Select a namespace}, choose your kiz username. Then click
  \textbf{Fork project} at the bottom.
\end{tasks}
\end{region}

\begin{region}[pitch=14.25]{340.3}
\begin{tasks}[start=5]
\item You are now on the project page of your own fork. Click \textbf{Code}
  there and copy the link under \textbf{Clone with HTTPS}.\\
  The link has the form: {\color{aslabblue}https://gitlab.uni-ulm.de/\textbf{USERNAME}/asp-ulm}
\end{tasks}
\end{region}

\begin{region}[pitch=13.2]{382.1}
\begin{tasks}[start=6]
\item As usual, we work with two terminals side by side -- the left one for
  editing, the right one for running commands.
\end{tasks}
\end{region}

\begin{region}{421.4}
\subheading{Left terminal (editing)}
\end{region}

\begin{region}[pitch=13.35]{435.7}
\begin{alphatasks}
\item \textbf{Clone your Git repository}\\
  Clone your fork with:\\
  \code{git clone https://gitlab.uni-ulm.de/USERNAME/asp-ulm}\\
  Then change into the directory with\\
  \code{cd asp-ulm}
\end{alphatasks}
\end{region}

\begin{region}[pitch=13.6]{503.0}
\begin{alphatasks}[start=2]
\item \textbf{Change a file as a test}\\
  Open the file \code{README.md} and write in a new line:
  \code{Here we go!}
\end{alphatasks}
\end{region}

\begin{region}{542.8}
\subheading{Right terminal (running commands)}
\end{region}

\begin{region}[pitch=12.65]{555.1}
\begin{alphatasks}
\item Also change into the project directory with \code{cd asp-ulm}\\
  Now try to commit the change with \code{git commit -a}\\
  At first, an error message appears, something like:
  \textbf{"Please tell me who you are…"}
\end{alphatasks}
\end{region}

\begin{region}[pitch=12.1]{604.5}
\begin{alphatasks}[start=2]
\item Then run the following two commands (with your own name and email
  address):\\
  \code{git config user.email "your.name@uni-ulm.de"}\\
  \code{git config user.name "Your Name"}
\end{alphatasks}
\end{region}

\begin{region}[pitch=12.2]{640.5}
\begin{alphatasks}[start=3]
\item Now run the commit again:\\
  \code{git commit -a}\\
  nano opens. Enter a short message, e.g.: \code{My first commit}\\
  Exit nano with Control + X and confirm with y.
\end{alphatasks}
\end{region}

\begin{region}[pitch=12.25]{689.1}
\begin{alphatasks}[start=4]
\item Transfer your commit to the server:\\
  \code{git push}\\
  You will be asked for your kiz username and password.\\
  Note: \textbf{the password is not displayed while you type}, not even as
  asterisks -- that is normal. Just type it blind and press Enter.
\end{alphatasks}
\end{region}

\begin{region}{750.1}
\begin{alphatasks}[start=5]
\item Switch back to the browser, reload the page, and check whether your
  change has been applied.
\end{alphatasks}
\end{region}

\nextpage
%% ===================================================================== 2 ==
\begin{region}{69.4}
\heading{Step 2: Programming the microcontroller with the Arduino framework}
\end{region}

\begin{region}{98.0}
\textbf{Preparation:} In both terminals, change into the subdirectory
\code{controller}.
\end{region}

\begin{region}[pitch=13]{124.0}
\textbf{Left terminal:} In the left terminal, open the file
\code{controller.ino} with \code{nano}. It contains the following program:
\end{region}

\begin{region}[pitch=10]{160.7}
\begin{asmlisting}[9]
  \aN{1} \aI{void} \aR{setup}()
  \aN{2} \{
  \aN{3}     \aR{Serial}.\aR{begin}(\aK{9600});
  \aN{4}     \aR{Serial}.\aR{println}(\aK{"Hello! Lots of A's coming up"});
  \aN{5} \}
  \aN{6}
  \aN{7} \aI{void} \aR{loop}()
  \aN{8} \{
  \aN{9}     \aR{Serial}.\aR{print}(\aK{"A"});
 \aN{10}     \aR{delay}(\aK{1000});
 \aN{11} \}
\end{asmlisting}
\end{region}

\begin{region}{285.0}
\subheading{Right terminal}
\end{region}

\begin{region}[pitch=12.15]{298.0}
\begin{numlist}
\item \textbf{Connect the microcontroller.}\\
  Attention: under WSL, you have to run the command \code{usb_check} every
  time you connect a microcontroller so that the USB port becomes visible in
  WSL.
\end{numlist}
\end{region}

\begin{region}{335.3}
\begin{numlist}[start=2]
\item \textbf{Compile and upload the program:} \code{make upload}
\end{numlist}
\end{region}

\begin{region}[pitch=12.33]{348.3}
\begin{numlist}[start=3]
\item \textbf{Start the serial monitor:} \code{make monitor}\\
  First, a line should appear with: \code{Hello! Lots of A’s coming up}
  and then -- as promised -- lots of \code{A}s. :-)\\
  You can exit the serial monitor with \code{Control-X}.
\end{numlist}
\end{region}

\begin{region}{410.2}
\subheading{Brief explanation of what happens}
\end{region}

\begin{region}[pitch=12.3]{422.6}
\code{controller.ino} actually contains only a small excerpt of the complete
C program that is compiled in the end. The functions \code{setup()} and
\code{loop()} are -- as shown here -- embedded internally in a
\code{main()} function:
\end{region}

\begin{region}[pitch=10]{469.9}
\begin{asmlisting}[9]
\aC{// Simplified representation – in reality, there are many more includes}
\aC{// and initializations here}
\mbox{}
\aI{int} main()
\{
    \aC{// setup() is called exactly once}
    setup();
\mbox{}
    \aC{// After that, loop() is called again and again}
    \aN{while} (\aK{1}) \{
        \aC{// With Arduino, an internal function checks, e.g., whether serial}
        \aC{// data is available:}
        \aN{if} (Serial.available()) \{ \aC{/* ... */}\}
\mbox{}
        loop();
    \}
\}
\end{asmlisting}
\end{region}

\begin{region}[pitch=12.2]{653.5}
So the function \code{setup()} is called exactly once -- in our example, to
initialize the serial interface and to output a line of text.\\
The function \code{loop()}, on the other hand, runs in an endless loop: in
each pass, the character \code{A} is output, followed by a wait of 1000
milliseconds (i.e., one second).
\end{region}

\nextpage
%% ===================================================================== 3 ==
\begin{region}{69.4}
\heading{What you should remember for today:}
\end{region}

\begin{region}[pitch=12.6]{81.2}
\begin{dashlist}
\item With \code{Serial.begin(9600)}, you initialize the serial interface --
  here with a baud rate of 9600.
\end{dashlist}
\end{region}

\begin{region}[pitch=12.6]{105.5}
\begin{dashlist}
\item With \code{Serial.print("foo")}, the string \code{"foo"} is output
  without an additional line break.
\end{dashlist}
\end{region}

\begin{region}[pitch=12.9]{129.8}
\begin{dashlist}
\item With \code{Serial.println("foo")}, the string \code{"foo"} is output
  followed by a line break (more precisely: \code{\r\n}).
\end{dashlist}
\end{region}

\begin{region}{167.7}
\heading{Open questions: How exactly does the serial interface work?}
\end{region}

\begin{region}[pitch=12.38]{180.0}
When \code{Serial.print("A");} is executed, you can observe on pin~1 (TX)
with the oscilloscope how the character \code{'A'} is transmitted -- as a
sequence of voltage changes that correspond to the bit pattern of
\code{0x41}.\\
The letter \code{'A'} has a so-called ASCII value and is represented as an
8-bit number: \code{01000001}. These bits are sent one after the other --
similar to Morse code -- by setting the TX pin \textbf{HIGH} or
\textbf{LOW}. At a baud rate of \textbf{9600}, each bit takes about
104\,µs.
\end{region}

\begin{region}[pitch=13]{267.2}
\textbf{But:} The data is not only output via \textbf{pin~1 (TX)}, but is
\textbf{also transmitted over the USB cable} -- with the help of a
\textbf{serial-to-USB converter} (often integrated into the microcontroller
or present on the board). That is why you can also see the transmitted
character in the \textbf{serial monitor on the computer}, without any
oscilloscope.
\end{region}

\begin{region}{331.2}
\heading{Finally}
\end{region}

\begin{region}[pitch=12.07]{343.3}
Change line~4 so that only \tcode{Hello!} followed by a line break is
output. Commit the changes and upload them to the repository with the
following command:\\
\code{git commit -a -m "Step 2"}\\
\code{git push}
\end{region}

\begin{region}{404.5}
\heading{Step 3: Playing around with the Arduino framework}
\end{region}

\begin{region}{419.4}
Modify the file \code{controller.ino} as follows:
\end{region}

\begin{region}[pitch=10]{430.1}
\begin{asmlisting}[9]
  \aN{1} \aI{void} \aR{setup}()
  \aN{2} \{
  \aN{3}     \aR{Serial}.\aR{begin}(\aK{9600});
  \aN{4} \}
  \aN{5}
  \aN{6} \aI{int} valueA0;
  \aN{7}
  \aN{8} \aI{void} \aR{loop}()
  \aN{9} \{
 \aN{10}     \aR{Serial}.\aR{print}(\aK{"X "});
 \aN{11}     \aR{Serial}.\aR{println}(valueA0);
 \aN{12}     \aR{delay}(\aK{100});
 \aN{13}     valueA0 += \aK{1};
 \aN{14} \}
\end{asmlisting}
\end{region}

\begin{region}{584.4}
\subheading{Explanation of the changes}
\end{region}

\begin{region}[free]{596.1}
\begin{dashlist}
\item In the function \code{setup()}, only the serial interface is
  initialized now.
\end{dashlist}
\end{region}

\begin{region}[pitch=12.45]{608.4}
\begin{dashlist}
\item In line~6, a global variable \code{valueA0} of type \code{int} is
  declared. Since it is defined outside a function, it is a global
  variable. In C/C++, global variables without an initial value
  automatically get the value \code{0}.
\end{dashlist}
\end{region}

\begin{region}[pitch=12.7,free]{645.0}
\begin{dashlist}
\item In line~10, the letter \code{'X'} followed by a space is output. This
  space will be important later.
\end{dashlist}
\end{region}

\begin{region}[pitch=12.9]{669.3}
\begin{dashlist}
\item In line~11, the current value of the variable \code{valueA0} is
  output, followed by a line break via \code{println}. This line break will
  also play a role later.
\end{dashlist}
\end{region}

\begin{region}[free]{693.9}
\begin{dashlist}
\item In line~13, \code{valueA0} is incremented by one in each pass of the
  \code{loop()} function.
\end{dashlist}
\end{region}

\begin{region}[pitch=12.1]{719.5}
\textbf{Task:} Compile the modified program, flash the microcontroller, and
compare the output in the serial monitor with the expected output.
\end{region}

\begin{region}{744.5}
\textbf{Finally:} \code{git commit -a -m "Step 3"} and \code{git push}
\end{region}

\nextpage
%% ===================================================================== 4 ==
\begin{region}{69.4}
\heading{Step 4: Tinkering and even more playing around}
\end{region}

\begin{region}[pitch=13]{96.7}
Grab a \textbf{breadboard, jumper wires, wires, a 10\,kΩ resistor, a push
button, and a potentiometer}, and build the following circuit:
\end{region}

\placegraphic{56.7}{112.4}{390.2}{breadboard-circuit.png}

\begin{region}{318.0}
\subheading{Important:}
\end{region}

\begin{region}{329.5}
\begin{dashlist}
\item The middle leg of the potentiometer is connected to pin A0.
\item One leg of the push button is connected to pin~2 (the other one via
  the resistor to GND).
\end{dashlist}
\end{region}

\begin{region}{366.3}
Edit the file \code{controller.ino} as follows:
\end{region}

\begin{region}[pitch=10]{389.0}
\begin{asmlisting}[9]
  \aN{1} \aI{void} \aR{setup}()
  \aN{2} \{
  \aN{3}     \aR{Serial}.\aR{begin}(\aK{9600});
  \aN{4}     \aR{pinMode}(\aK{2}, \aK{INPUT});
  \aN{5} \}
  \aN{6}
  \aN{7} \aI{int} valueA0;
  \aN{8}
  \aN{9} \aI{void} \aR{loop}()
 \aN{10} \{
 \aN{11}     \aN{if} (\aR{digitalRead}(\aK{2})) \{
 \aN{12}         \aR{Serial}.\aR{println}(\aK{"A"});
 \aN{13}     \}
 \aN{14}
 \aN{15}     \aI{int} analogValue = \aR{analogRead}(A0);
 \aN{16}     \aN{if} (valueA0 != analogValue) \{
 \aN{17}         valueA0 = analogValue;
 \aN{18}         \aR{Serial}.\aR{print}(\aK{"X "});
 \aN{19}         \aR{Serial}.\aR{println}(valueA0);
 \aN{20}     \}
 \aN{21}     \aR{delay}(\aK{100});
 \aN{22} \}
\end{asmlisting}
\end{region}

\begin{region}{621.3}
\subheading{Task}
\end{region}

\begin{region}{633.3}
Before the program is explained in detail, first observe for yourselves what
it does:
\end{region}

\begin{region}{644.8}
\begin{dashlist}
\item Compile the program, flash the microcontroller, and start the serial
  monitor.
\end{dashlist}
\end{region}

\begin{region}[pitch=12.5]{656.8}
\begin{dashlist}
\item Observe what happens when you turn the potentiometer -- especially at
  its extreme positions.
\end{dashlist}
\end{region}

\begin{region}{680.8}
\begin{dashlist}
\item Observe what happens when you press the push button.
\end{dashlist}
\end{region}

\begin{region}{706.3}
\textbf{Finally:} \code{git commit -a -m "Step 4"} and \code{git push}
\end{region}

\nextpage
%% ===================================================================== 5 ==
\begin{region}{69.4}
\subheading{Explanation of the changes}
\end{region}

\begin{region}{81.5}
\begin{dashlist}
\item In \textbf{line~4}, \code{pinMode(2, INPUT)} configures digital pin~2
  as an input.
\end{dashlist}
\end{region}

\begin{region}[pitch=12.72]{94.5}
\begin{dashlist}
\item In \textbf{lines 11 to 13}, you see an \code{if} statement. The
  condition checks the value returned by \code{digitalRead(2)}. This function
  returns \code{1} (corresponding to \textbf{HIGH} or \code{true}) if a
  voltage of around 5\,V is present at pin~2, and \code{0} (corresponding to
  \textbf{LOW} or \code{false}) if there is almost no voltage.\\
  If the return value is 1, the character \code{'A'} is sent via the serial
  interface in line~12 -- including a line break via
  \code{Serial.println("A")}
\end{dashlist}
\end{region}

\begin{region}[pitch=12.53,free]{170.2}
\begin{dashlist}
\item In \textbf{line~15}, a local variable \code{analogValue} of type
  \code{int} is declared and initialized with the return value of
  \code{analogRead(A0)}. Pin~A0 is an analog input. A voltage between 0\,V
  and 5\,V is converted by the microcontroller into a digital value between
  \code{0} and \code{1023} (i.e., 10-bit resolution).
\end{dashlist}
\end{region}

\begin{region}[pitch=12.65]{219.8}
\begin{dashlist}
\item The \code{if}\textbf{ statement in lines 16 to 20} checks whether the
  new analog value \code{analogValue} differs from the previously stored
  value \code{valueA0}. Only if the value has changed,
\end{dashlist}
\end{region}

\begin{region}[x=68.7,right=0]{257.1}
\begin{dashlist}
\item is the new value stored in \code{valueA0} in \textbf{line~17},
\end{dashlist}
\end{region}

\begin{region}[x=68.7,right=0]{270.1}
\begin{dashlist}
\item is an output generated via the serial interface in
  \textbf{lines 18--19} (in the format \code{X <value>}).
\end{dashlist}
\end{region}

\begin{region}[x=68.7,right=0]{283.1}
If the value remains unchanged, no output is generated. This avoids
unnecessary repetition.
\end{region}

\begin{region}{320.1}
\heading{What you should remember for today:}
\end{region}

\begin{region}[pitch=12.3,free]{331.8}
\begin{dashlist}
\item With \code{pinMode(2, INPUT)}, a digital pin is configured as an
  input.
\item \code{digitalRead(2)} returns \code{1} if a voltage is present at the
  pin, and \code{0} if not.
\item With \code{analogRead(A0)}, a voltage from 0\,V to 5\,V is converted
  into a value from \code{0} to \code{1023}.
\end{dashlist}
\end{region}

\begin{region}{406.0}
\heading{Step 5: Fine-tuning}
\end{region}

\begin{region}{430.0}
The behavior of the microcontroller is to be refined:
\end{region}

\begin{region}[pitch=13.03]{442.0}
\begin{dashlist}
\item \textbf{When the push button is pressed,} an \code{'A'} is to be sent
  via the serial interface.
\item \textbf{As long as the push button remains pressed,} no further
  character may be sent.
\item \textbf{When the push button is released,} an \code{'a'} is to be
  sent.
\item \textbf{As long as the push button is not pressed,} no further
  \code{'a'} may be sent.
\end{dashlist}
\end{region}

\begin{region}[pitch=12.1]{494.3}
In short: the characters \code{'A'} and \code{'a'} should only be sent on a
change of state -- not while the button is being held.
\end{region}

\begin{region}{531.3}
\subheading{Task}
\end{region}

\begin{region}[pitch=13]{555.6}
\begin{numlist}
\item Create a global variable \code{pressedA} of type \code{bool}. For it
  to be global, it must be declared \textbf{outside a function} -- e.g.,
  after line~7:
\end{numlist}
\end{region}

\begin{region}[x=74.7,right=0]{580.7}
\begin{asmlisting}[9]
\aI{bool} pressedA;
\end{asmlisting}
\end{region}

\begin{region}[x=74.7,right=0,pitch=12.3]{592.9}
Since global variables in C/C++ are automatically initialized with the
value \code{0}, \code{pressedA} is initially \code{false}.
\end{region}

\begin{region}[free]{617.5}
\begin{numlist}[start=2]
\item Change the \code{if} statement in \code{loop()} to
\end{numlist}
\end{region}

\begin{region}[x=74.7,right=0,pitch=10]{640.3}
\begin{asmlisting}[9]
\aN{if} (!pressedA \&\& \aR{digitalRead}(\aK{2})) \{
    \aR{Serial}.\aR{println}(\aK{"A"});
    pressedA = \aK{true};
\}
\end{asmlisting}
\end{region}

\begin{region}[x=74.7,right=0,pitch=12.3]{681.9}
The \code{!} negates the logical value of \code{pressedA}, \code{&&} is the
logical AND. With this condition, an \code{'A'} is only output if the push
button was not pressed before.\\
At the same time, we remember with \code{pressedA = true} that the push
button is now pressed.
\end{region}

\begin{region}[pitch=12.1]{718.5}
\begin{numlist}[start=3]
\item Compile the program, flash the microcontroller, and observe the output
  in the serial monitor. With this change, part of the desired behavior is
  already achieved. However, the output of \code{'a'} on release is still
  missing. With a second \code{if} statement and a bit of puzzling, you
  should be able to reach the goal completely.
\end{numlist}
\end{region}

\begin{region}{767.8}
\begin{numlist}[start=4]
\item \textbf{If everything works:} \code{git commit -a -m "Step 5"} and
  \code{git push}
\end{numlist}
\end{region}

\nextpage
%% ===================================================================== 6 ==
\begin{region}{69.4}
\heading{Step 6: From the hello world program to the serial interface}
\end{region}

\begin{region}[pitch=12.24]{94.8}
In both terminals, change into the directory \code{python}. This directory
is located next to the current directory \code{controller} -- you get there
with:\\
\code{cd ..}\\
\code{cd python}\\
or in one go:\\
\code{cd ../python}
\end{region}

\begin{region}{180.3}
There you will find the file \code{serial-reader.py} with the following
content:
\end{region}

\begin{region}{203.0}
\begin{asmlisting}[9]
\aR{print}(\aK{"Hello World"})
\end{asmlisting}
\end{region}

\begin{region}{226.3}
This is still a classic hello world program. Run it with:
\end{region}

\begin{region}{250.6}
\code{python serial-reader.py}
\end{region}

\begin{region}{274.9}
The output should -- not very surprisingly -- simply be: \code{Hello World}
(with a line break).
\end{region}

\begin{region}{299.0}
Now we dare to take a giant step forward: we modify the Python program so
that it reads the serial interface!
\end{region}

\begin{region}{335.0}
Along the way, we also learn (learning by doing) a few basic language
features of Python.
\end{region}

\begin{region}{371.2}
Modify the program \code{serial-reader.py} as follows:
\end{region}

\begin{region}[pitch=10]{394.0}
\begin{asmlisting}[9]
  \aN{1} \aM{import} serial
  \aN{2} \aM{import} time
  \aN{3}
  \aN{4} ser = serial.Serial(\aK{"/dev/ttyUSB0"}, \aK{9600})
  \aN{5}
  \aN{6} \aR{print}(\aK{"Waiting for serial input..."})
  \aN{7}
  \aN{8 while} \aR{True}:
  \aN{9}     \aN{while} ser.in\_waiting > \aK{0}:
 \aN{10}         ch = ser.read(\aK{1})
 \aN{11}         \aR{print}(ch)
 \aN{12}     time.sleep(\aK{0.1})
\end{asmlisting}
\end{region}

\begin{region}{528.2}
\subheading{Indentation is not negotiable}
\end{region}

\begin{region}{553.2}
When it comes to wrong indentation, Python is \textbf{just as tolerant as I
am} -- that is, \textbf{not at all}.
\end{region}

\begin{region}[pitch=12.53]{578.2}
Unlike many other programming languages, Python has \textbf{no curly
braces} to delimit loop or function blocks. Instead, the Python interpreter
-- just like the programmer -- recognizes \textbf{solely from the
indentation} which statements belong to which block.\\
In our example:
\end{region}

\begin{region}{640.3}
\begin{dashlist}
\item The outer \code{while True:} loop starts in \textbf{line 8} and
  comprises \textbf{lines 9 to 12}.
\end{dashlist}
\end{region}

\begin{region}[pitch=13.9]{653.3}
\begin{dashlist}
\item The inner loop \code{while ser.in_waiting > 0:} starts in
  \textbf{line 9} and comprises only \textbf{lines 10 and 11}.
\end{dashlist}
\end{region}

\begin{region}[pitch=13]{678.9}
\begin{dashlist}
\item \code{time.sleep(0.1)} no longer belongs to the inner loop -- you can
  tell this solely from the fact that it is indented at the same level as the
  inner \code{while}.
\end{dashlist}
\end{region}

\nextpage
%% ===================================================================== 7 ==
\begin{region}{69.4}
\subheading{Task}
\end{region}

\begin{region}{93.8}
Run the program with \code{serial-reader.py}
\end{region}

\begin{region}{117.8}
You should see output that looks roughly like this (if necessary, turn the
potentiometer or press the push button):
\end{region}

\begin{region}{154.1}
\code{b'X'}\\
\code{b' '}\\
\code{b'4'}\\
\code{b'0'}\\
\code{b'4'}\\
\code{b'\r'}\\
\code{b'\n'}
\end{region}

\begin{region}[pitch=13]{250.4}
Exit the program with \code{Control-C}.\\
Update your repository with \code{git commit -a -m "Step 6"} and
\code{git push}
\end{region}

\begin{region}{288.4}
\subheading{Explanation of the output}
\end{region}

\begin{region}{311.9}
\begin{dashlist}
\item The program outputs one received byte per line.
\item The data comes directly from the serial interface of the
  microcontroller.
\end{dashlist}
\end{region}

\begin{region}[pitch=12.8]{335.9}
\begin{dashlist}
\item In Python, such bytes are represented as so-called byte objects --
  you can recognize them by the prefix \code{b'...'}.
\end{dashlist}
\end{region}

\begin{region}[pitch=12.5]{360.2}
\begin{dashlist}
\item Unlike normal character strings (strings), a byte object here contains
  exactly one byte, for example:
\end{dashlist}
\end{region}

\begin{region}[x=68.7,right=0,pitch=12.3]{384.4}
\begin{dashlist}
\item \code{b'X'} stands for the byte with the ASCII value \code{0x58}, i.e.,
  the letter \code{X}
\item \code{b' '} is the space character (\code{0x20})
\item \code{b'\r'} is the control character carriage return (\code{0x0D})
\item \code{b'\n'} is the line feed (\code{0x0A})
\end{dashlist}
\end{region}

\begin{region}[x=68.6,right=0,pitch=12.3]{434.2}
These last two characters are produced when \code{Serial.println(...)} is
used in the Arduino program -- because \code{println} automatically sends
\code{\r\n} at the end.
\end{region}

\begin{region}{471.5}
\subheading{Technical background}
\end{region}

\begin{region}[pitch=12.1]{495.8}
The variable \code{ser}, which is created by \code{serial.Serial(...)},
represents a special data structure that manages access to the serial
interface.
\end{region}

\begin{region}[pitch=12.1]{520.8}
An internal \textbf{buffer} is also created, in which received bytes are
stored temporarily.
\end{region}

\begin{region}[x=68.7,right=0,pitch=12.5]{544.4}
\begin{dashlist}
\item Whenever a byte is received via the serial interface, it first ends
  up in this buffer.
\end{dashlist}
\end{region}

\begin{region}[x=68.7,right=0,pitch=12.6]{568.6}
\begin{dashlist}
\item The variable \code{ser.in_waiting} always contains the number of bytes
  currently stored in the buffer.
\end{dashlist}
\end{region}

\begin{region}[x=68.7,right=0]{593.2}
\begin{dashlist}
\item With \code{ser.read(1)}, you can read \textbf{the oldest byte} from the
  buffer.
\end{dashlist}
\end{region}

\begin{region}[pitch=12.1]{607.2}
This buffer works according to the principle of a \textbf{FIFO data
structure} \textit{(First In, First Out)}:\\
The byte received first is also the first one to be read.
\end{region}

\begin{region}{644.2}
\subheading{What you should remember for today:}
\end{region}

\begin{region}[x=68.7,right=0]{668.9}
\begin{dashlist}
\item With \code{import serial}, the module pyserial is included -- it must
  be installed.
\end{dashlist}
\end{region}

\begin{region}[x=68.7,right=0,pitch=12.45,free]{681.2}
\begin{dashlist}
\item \code{ser = serial.Serial(...)} opens the serial interface (here:
  \code{/dev/ttyUSB0}) with the specified baud rate.\\
  All further communication with the interface takes place via the
  variable \code{ser}.
\end{dashlist}
\end{region}

\begin{region}[x=68.7,right=0,pitch=12.33,free]{717.8}
\begin{dashlist}
\item \code{ser.read(1)} reads one byte from the interface and returns it as
  a byte object.
\item The condition \code{ser.in_waiting > 0} checks whether data has
  already been received.
\item The prefix \code{b'...'} means in Python: this is a byte object, not a
  string.
\item \code{import time} loads the built-in module \code{time} (it does not
  need to be installed separately).
\end{dashlist}
\end{region}

\begin{region}[x=68.7,right=0,pitch=12.6,free]{767.1}
\begin{dashlist}
\item With \code{time.sleep(0.1)}, a pause of a tenth of a second is
  inserted -- this relieves the processor and prevents unnecessary loop
  activity.
\end{dashlist}
\end{region}

\nextpage
%% ===================================================================== 8 ==
\begin{region}{69.4}
\heading{Step 7: Homework!}
\end{region}

\begin{region}{93.5}
You can program really great things with Python -- you just need to get
familiar with a few basics. To get you fit for this, here is a small
homework assignment:\\
At the university, we have licenses for an interactive online tutorial:
\end{region}

\begin{region}{133.9}
\makebox[18.1bp][l]{\emoji{1F449}}\bluelink{https://www.uni-ulm.de/nawi/physik/fuer-studierende/angebote-fuer-studierende/opb-online-python-brueckenkurs/}{Online Python Bridging Course (OPB)}
\end{region}

\placegraphic{56.7}{138.5}{132.1}{qr-code-opb.png}

\begin{region}{283.5}
\subheading{Task until the next session (in two weeks):}
\end{region}

\begin{region}[pitch=13]{295.5}
\begin{dashlist}
\item Work through \textbf{chapters 1 to 11}
\item Pass the \textbf{challenges ``Chapters 1--7''} and
  \textbf{``Chapters 8--11''}
\end{dashlist}
\end{region}

\begin{region}{334.5}
The chapters are \textbf{short and doable} -- about 5--10 minutes each. So
no panic!
\end{region}

\begin{region}{358.5}
Afterwards, you will have a good understanding of central language elements
of Python, e.g.:
\end{region}

\begin{region}[pitch=13]{370.5}
\begin{dashlist}
\item \textbf{Variables} (chapter~4)
\item \textbf{Numbers} (chapter~5)
\item \textbf{Strings} (chapter~6)
\item \textbf{Booleans} (chapter~7)
\item \textbf{If statements} (chapter~8)
\item \textbf{Lists} (chapter~9)
\item \textbf{For loops} (chapter~10)
\item \textbf{While loops} (chapter~11)
\end{dashlist}
\end{region}

\begin{region}[pitch=16.4]{486.5}
If reading the serial interface was not motivation enough for you -- just
wait for the next chill-out step\,\emoji{1F609}
\end{region}

\begin{region}{533.2}
\emoji{1F3AE} \textbf{Chill-out step: Breakout -- almost done…}
\end{region}

\begin{region}[pitch=14.2]{559.4}
In the directory \code{python}, you will also find the Python program
\code{breakout.py}. It is not yet the finished game -- but exactly what you
will be able to program yourselves next time, once you have diligently
worked through the online Python course. \emoji{1F609}
\end{region}

\begin{region}{614.0}
But now: just try it out! Run the program: \code{python breakout.py}
\end{region}

\begin{region}{639.0}
\subheading{But afterwards, we clean up}
Remove \code{breakout.py} from your repository -- because writing it
yourselves is more fun!
\end{region}

\begin{region}{675.6}
\code{git rm breakout.py}\\
\code{git commit -a -m "Done!"}\\
\code{git push}
\end{region}

\begin{region}{723.7}
Next time, you will build the game yourselves step by step --\\
and understand exactly how Python talks to hardware.
\end{region}

\end{document}
